\section{Matrices over Euclidean Domains}

At this point of the course, we are interested in proving and generalising the previous result on finite/finitely generated abelian groups -- which are just \(\ZZ\)-modules. Furthermore, a matrix \(\vm{A}\) on \(\ZZ\) (or any Euclidean domain) gives a module homomorphism \(\functiondomain{\vm{A}}{\ZZ^m}{\ZZ^n}\) (or generally, \(\functiondomain{\vm{A}}{R^m}{R^n}\)). The groups/modules we are interested in will be \(\ZZ^n / \img \vm{A}\), since they are finitely generated, and this perspective will give us a classification.

Now, we fix a Euclidean domain \(R\), and the associated Euclidean function \(\functiondomain{\phi}{R \setminus \brce{0}}{\NN_0}\). By Euclid's algorithm, for \(a, b \in R\), we can find \(x, y \in R\), such that
\[
    ax + by = \gcd\brce{a, b}.
\]

\setcounter{MaxMatrixCols}{20}

\begin{definition}[Elementary Row/Column Operations]
    \label{def:elementary-row-column-operation}%
    Let \(\vm{A}\) be a matrix with entries in \(R\). The \emph{elementary row operations} are the following:
    \begin{enumerate}
        \item Add \(c \in R\) times the \(i\)th row to the \(j\)th row, i.e., left-multiply by the matrix
        \[
            \begin{pmatrix}
                1& & & &0\\
                 &1& & & \\
                 & &\ddots&c& \\
                 & & &1& \\
                0& & & &1
            \end{pmatrix}
        \]
        where \(c\) is in the \(\para{j, i}\)-th position.

        \item Swap rows \(i\) and \(j\), i.e., left-multiply by the matrix
        \[
            \begin{pmatrix}
                1& & & & & & & & & &0\\
                 & \ddots & & & & & & & & & \\
                 & &1& & & & & & & & \\
                 & & &0&0&\cdots&0&1& & & \\
                 & & &0&1& &0&0& & & \\
                 & & &\vdots& & \ddots & & \vdots & & & \\
                 & & &0&0& &1&0& & & \\
                 & & &1&0&\cdots&0&0& & & \\
                 & & & & & & & &1& & \\
                 & & & & & & & & & \ddots & \\
                0& & & & & & & & & &1
            \end{pmatrix}
        \]
        where the \(\para{i, j}\) and \(\para{j, i}\) entries are \(1\), and all diagonal entries are \(1\) apart from the entries \(\para{i, i}\) and \(\para{j, j}\) are \(0\).

        \item Multiply row \(i\) by a \emph{unit} \(c \in R\multiplicative\), i.e., left-multiply by the matrix
        \[
            \begin{pmatrix}
                1& & & & & &0\\
                 &\ddots& & & & &\\
                 & &1& & & & \\
                 & & &c& & & \\
                 & & & &1& & \\
                 & & & & &\ddots& \\
                0& & & & & &1  
            \end{pmatrix} = \diag\para{1, \ldots, 1, c, 1, \ldots, 1}
        \]
        where \(c\) is in the \(\para{i, i}\)th position.
    \end{enumerate}

    \emph{Elementary column operations} are defined similarly.
\end{definition}

\begin{definition}[Equivalent Matrices]
    \label{def:equivalent-matrix}%
    Matrices \(\vm{A}\) and \(\vm{B}\) are \emph{equivalent}, if \(A\) can be obtained from \(\vm{B}\) via row/column operations. Equivalently,
    \[
        \vm{B} = \vm{Q} \vm{A} \vm{T}\inverse
    \]
    for \(\vm{Q}\) and \(\vm{T}\) invertible.
\end{definition}

Observe that the matrix
\[
    \begin{pmatrix}
        2 & 0 \\ 0 & 0
    \end{pmatrix}
\]
cannot be further reduced in \(\ZZ\) -- unlike over a field, where a matrix can always be reduced to a number of \(1\)s followed by \(0\)s in the diagonal.

\begin{theorem}[Smith Normal Form]
    \label{thm:smith-normal-form}%
    Any \(m \times n\) matrix is equivalent to one of the form
    \[
        \begin{pmatrix}
            d_1 & & & & & 0 \\
            & \ddots & & & & & \\
            & & d_r & & & \\
            & & & 0 & & \\
            & & & & \ddots & \\
            0 & & & & & 0
        \end{pmatrix} = \diag\para{d_1, \ldots, d_r, 0, \ldots, 0}
    \]
    where \(d_i \divides d_{i + 1}\). The \(d_i\)s are called the \emph{invariant factor} of the matrix.
\end{theorem}

\begin{proof}
    Let the matrix be \(\vm{A}\). If \(\vm{A}\) is \(\zeroM\), then we are done.

    So assume \(A_{i, j} \neq 0\) for some \(i, j\). Then perform row/column operations, so that \(A_{1, 1}\) is non-zero.

    \begin{enumerate}
        \item If \(A_{1, j}\) is not divisible by \(A_{1, 1}\), then we write \(A_{1, j} = q A_{1, 1} + r\) for some \(q, r \in R\), and \(\phi\para{r} < \phi\para{A_{i, i}}\). By column operations, make \(\para{i, j}\) entry equal to \(r\). Now swap so that the \(\phi\) value of \(\para{1, 1}\) reduces.

        \item Similarly, we do the same for \(A_{j, 1}\) using row operations. After finitely many such moves, \(A_{1, 1}\) must divide all the first row and the first column.

        By row and column operations again, we can make the first row and column zero outside the \(\para{1, 1}\) entry. This gives the matrix
        \[
            \begin{pmatrix}
                d & 0 & \cdots & 0 \\
                0 & & & \\
                \vdots & & \vm{C} & \\
                0 & & &
            \end{pmatrix}
        \]

        \item Now, suppose \(A_{i, j}\) in \(\vm{C}\) is not divisible by \(d\). Then we write \(A_{i, j} = dq + r\) with \(\phi\para{r}, \phi\para{d}\). We now add column \(1\) to column \(j\), and subtract \(q\) times row \(1\) from row \(i\). This gives \(r\) in the \(\para{i, j}\) entry. Now, swap back this \(r\) to the \(\para{1, 1}\) position, and then we have \(\phi\para{r} < \phi\para{d}\).

        \item We return to 1 and 2, and replace \(\vm{A}\) with the new matrix
        \[
            \begin{pmatrix}
                d' & 0 & \cdots & 0 \\
                0 & & & \\
                \vdots & & \vm{C}' & \\
                0 & & &
            \end{pmatrix}
        \]
        with \(\phi\para{d'} < \phi\para{r} < \phi\para{d}\). Repeat finitely many times, so eventually \(\vm{A}\) becomes
        \[
            \begin{pmatrix}
                d & 0 & \cdots & 0 \\
                0 & & & \\
                \vdots & & \vm{C} & \\
                0 & & &
            \end{pmatrix},
        \]
        but this time with \(d \divides C_{i, j}\) for all entries \(\para{i, j}\).

        \item Now, we recurse the algorithm, and the output is as claimed.
    \end{enumerate}
\end{proof}

\begin{definition}[Minors, Fitting Ideals]
    \label{def:minor-fit-ideals}%
    For a matrix \(\vm{A}\), a \emph{\(k \times k\) minor} is the determinant of a \(k \times k\) sub-matrix. The \emph{\(k\)th Fitting ideal}
    \[
        \Fit_k \para{\vm{A}} \idealsub R
    \]
    is the ideal generated by the \(k \times k\) minors of \(\vm{A}\).
\end{definition}

\begin{proposition}[Row/Column Operations preserve Fitting Ideals]
    \label{prop:row-column-operations-preserve-fit}%
    If \(\vm{A}\) and \(\vm{B}\) are equivalent matrices, then
    \[
        \Fit_k \para{\vm{A}} = \Fit_k \para{\vm{B}}
    \]
    for all \(k\).
\end{proposition}

\begin{proof}
    It suffices to check that individual row/column operations do not change the Fitting ideals. Since \(\Fit_k\para{\vm{A}} = \Fit_k\para{\vm{A}\transpose}\), it suffices to check this for the row operations. Swapping row are straightforward.

    Consider \(\vm{C}\) a \(k \times k\) sub-matrix, and consider adding \(c\) times row \(i\) of \(\vm{A}\) to row \(j\). This gives a new sub-matrix \(\vm{C}'\). We show that \(\det\para{\vm{C}'} \in \Fit\para{\vm{A}}\).

    If the \(j\)th row is outside \(\vm{C}\), then this does not change anything.

    If both rows lie in \(\vm{C}\), this follows from the properties of the determinant.

    Now, suppose the \(j\)th row is in \(\vm{C}\), and the \(i\)th row is not. Suppose row \(i\) is given by
    \[
        \begin{bmatrix}
            f_1 & \cdots & f_k
        \end{bmatrix}.
    \]

    Then the new row \(j\) is given by
    \[
        \begin{bmatrix}
            C_{j, 1} + c f_1 & \cdots & C_{j, k} + c f_k
        \end{bmatrix}.
    \]

    Now compute \(\det \para{\vm{C}'}\) along this row. Using the additivity of the determinant, we have
    \[
        \det\para{\vm{C}'} = \det\para{\vm{C}} + c \det\para{\vm{D}},
    \]
    where \(\vm{D}\) is obtained by row \(j\) of \(\vm{C}\) with the original row \(i\). But both \(\vm{C}\) and \(\vm{D}\) are in fact \(k \times k\) sub-matrices, so we are done.

    So we concluded
    \[
        \Fit_k \para{\vm{B}} \subseteq \Fit_k \para{\vm{A}}.
    \]

    But such operations are invertible, so
    \[
        \Fit_k \para{\vm{A}} = \Fit_k \para{\vm{B}}
    \]
    as desired.
\end{proof}

\begin{corollary}[Fitting Ideal from Smith Normal Form, Uniqueness of Invariant Factors]
    \label{cor:fit-ideal-snf-uniqueness-if}%
    If \(\vm{A}\) has Smith normal form given by
    \[
        \diag\para{d_1, \ldots, d_r, 0, \ldots, 0},
    \]
    then
    \[
        \Fit_k \para{\vm{A}} = \para{d_1 \cdots d_k},
    \]
    and in particular, the \(d_i\)s are unique, up to associates.
\end{corollary}

\begin{proposition}[Size of Generating Set of a Submodule of a Free Module]
    \label{prop:size-generating-set-free-module}%
    Let \(R\) be a PID. Then any submodule of \(R^m\) can be generated by \(m\) or fewer elements.
\end{proposition}

\begin{proof}
    We proceed by induction on \(m\). Let \(N \submodule \RR^m\) be a submodule. We consider the ideal
    \[
        I = \set{r \in R}{\para{r, r_2, \ldots, r_m} \in N \text{ for some }r_i}.
    \]

    Observe \(I \idealsub R\). Since \(R\) is a PID, \(I = \para{a}\) for some \(a \in R\). So let
    \[
        n = \para{a, a_2, \ldots, a_m} \in N.
    \]

    So we know for \(\para{r_1, \ldots, r_m} \in N\), \(a \divides r_1\), and hence
    \[
        \para{r_1, \ldots, r_m} - r \cdot n = \para{0, r_2 - r a_2, \ldots, r_m - r a_m} \in N.
    \]

    So everything in \(N\) is a multiple of \(n\), plus something in \(N \cap \para{\brce{0} \times R^{m - 1}}\). Now we apply induction.
\end{proof}

\begin{theorem}[Basis of Submodule may Extend to Basis of Free Module]
    \label{thm:basis-submodule-extend-space}%
    Let \(R\) be a Euclidean domain, and \(N \submodule R^m\). Then there exists a basis \(v_1, \ldots, v_m\) of \(N\), such that \(N\) is generated by \(d_1 v_1, \ldots, d_r v_r\) by some \(0 \leq r \leq m\), \(d_i \in R\), with \(d_i \divides d_{i + 1}\).
\end{theorem}

Note the multiples \(d_1, \ldots, d_r\) are the crucial difference between the Euclidean domain and the field case.

\begin{proof}
    By \autoref{prop:size-generating-set-free-module}, \(N\) can be generated by \(x_1, \ldots, x_m \in R^m\) with \(n \geq m\). Consider the matrix
    \[
        \begin{pmatrix}
            \vert & \vert &  & \vert \\
            x_1 & x_2 & \cdots & x_m \\
            \vert & \vert &  & \vert
        \end{pmatrix}.
    \]

    We may put the matrix into Smith normal form and obtain
    \[
        \diag\para{d_1, \ldots, d_r, 0, \ldots, 0}
    \]
    by row operations.

    Row operations change the basis for \(R^m\). Column operations change generators for \(N\).

    By that new basis, \(N\) is generated by \(\brce{d_i v_i}\), as required.
\end{proof}

\begin{theorem}[Submodule of Module over Euclidean Domain is Free]
    \label{thm:submodule-module-euclidean-domain-free}%
    For \(R\) a Euclidean domain, any submodule of \(R^m\) is free.
\end{theorem}

\begin{proof}
    In the proof of \autoref{thm:basis-submodule-extend-space}, the generating set \(\brce{d_i v_i}\) is linearly independent, because any linear combination of the \(d_i v_i\) which gives zero, also gives a linear combination for the \(v_i\) themselves, but \(v_i\) are free.
\end{proof}

\begin{remark}
    This in fact fails over other rings. The submodule \(\para{x, y} \idealsub \CC\brac{x, y}\) is not free.
\end{remark}

\begin{theorem}[Classification of Finitely-Generated Modules over Euclidean Domains]
    \label{thm:classification-finitely-generated-modules-euclidean-domains}%
    Let \(R\) be a Euclidean domain. Let \(M\) be a finitely-generated \(R\)-module. Then
    \[
        M \isomorphic R \oplus \cdots \oplus R \oplus R / \para{d_1} \oplus \cdots \oplus R / \para{d_r}
    \]
    where \(d_i \neq 0\), \(d_i \divides d_{i + 1}\).
\end{theorem}

\begin{proof}
    Let \(\functiondomain{\gamma}{R^m}{M}\) be a surjective \(R\)-module homomorphism. We know \(\ker \gamma \leq R^m\), and we may take a basis \(v_1, \ldots, v_m\) such that \(\ker \gamma\) is generated by \(d_1 v_1, \ldots, d_r v_r\) for some \(0 \leq r \leq m\), with \(d_i \divides d_{i + 1}\).

    Then by \autoref{thm:first-isomorphism-theorem-modules} (First Isomorphism Theorem for Modules), we have
    \[
        M \isomorphic R^m / \para{\para{d_1, 0, \ldots, 0}, \para{0, d_2, \ldots, 0}, \ldots, \para{0, \ldots, d_r, \ldots, 0}}.
    \]
\end{proof}

\begin{remark}
    When \(R = \ZZ\), this gives the Classification of Finitely-Generated Abelian Groups (see \autoref{thm:classification-finite-abelian-group}).
\end{remark}

\begin{example*}[Identification of Presentation of Finitely Generated Abelian Group]
    Let \(G\) be an abelian group generated \(a, b, c\), with \(2a + 3b + c = 0\), \(a + 2b = 0\), and \(5a + 6b + 7c = 0\).

    In other words, we have a surjective homomorphism \(\functiondomain{\gamma}{\ZZ^3}{G}\), sending \(\para{1, 0, 0}\) to \(a\), \(\para{0, 1, 0}\) to \(b\), and \(\para{0, 0, 1}\) to \(c\).

    The kernel of \(\gamma\) is therefore given by the image of the matrix
    \[
        \vm{A} = \begin{pmatrix}
            2 & 1 & 5 \\
            3 & 2 & 6 \\
            1 & 0 & 7
        \end{pmatrix},
    \]
    where we view this as a map \(\ZZ^3\) to \(\ZZ^3\). By construction,
    \[
        G \isomorphic \ZZ^3 / \img \para{\vm{A}}.
    \]

    To calculate the Smith Normal Form, we proceed by calculating Fitting ideals of 
    \(vm{A}\). We note
    \[
        \Fit_1 \para{\vm{A}} = \Fit_2 \para{\vm{A}} = \para{1},
    \]
    while
    \[
        \Fit_3 \para{\vm{A}} = \para{3}.
    \]

    Since \(d_1 = d_2 = 1, d_3 = 3\), we have
    \[
        G \isomorphic \ZZ / \para{3}.
    \]
\end{example*}

\begin{proposition}[Chinese Remainder Theorem for Modules]
    \label{prop:crt-modules}%
    Let \(R\) be a Euclidean domain, with \(a, b \in R\), such that \(\gcd\para{a, b} = 1\). Then
    \[
        R / \para{ab} \isomorphic R/\para{a} \oplus R/\para{b}.
    \]
\end{proposition}

\begin{proof}
    Consider the map
    \[
        \function{\phi}{R / \para{a} \oplus R / \para{b}}{R / \para{ab}}{\para{r_1 + \para{a}, r_2 + \para{b}}}{\para{b r_1 + a r_2 + \para{ab}}.}
    \]

    We may check that this is well-defined.

    \emph{Surjectivity} arises from the coprime condition, since we have the Euclidean algorithm on a Euclidean domain. This means \(1 + \para{ab}\) lies in the image of \(\phi\). This generates the whole module.

    \emph{Injectivity} arises from calculating the kernel. If \(b r_1 + a r_2 \in \para{ab}\). We write
    \[
        b r_1 + a r_2 = a b x
    \]
    for some \(x \in R\). Since \(a \divides a r_2\) and \(a \divides a b x\), we have \(a \divides b r_1\), and hence \(a \divides r_1\). Similarly, \(b \divides r_2\).

    Therefore, \(r_1 \equiv 0 \pmod a\) and \(r_2 \equiv 0 \pmod b\), which means the kernel is trivial.
\end{proof}

\begin{theorem}[Primary Decomposition of Classification]
    \label{thm:primary-decomposition-classification}%
    Let \(R\) be a Euclidean domain, and \(M\) a finitely-generated \(R\)-module. Then
    \[
        M \isomorphic N_1 \oplus \cdots N_t,
    \]
    where each \(N_i\) is either \(N_i \isomorphic R\), or \(N_i \isomorphic R / \para{{p_i}^{n_i}}\) for some prime \(p_i\) and exponent \(n_i \geq 1\).
\end{theorem}

\begin{proof}
    From the previous classification, we have
    \[
        M \isomorphic R \oplus \cdots \oplus R \oplus R / \para{d_1} \oplus \cdots \oplus R / \para{d_r}.
    \]

    Now we write each \(d_i\) as a product of primes in \(R\), and apply \autoref{prop:crt-modules} (Chinese Remainder Theorem for Modules).
\end{proof}